\documentclass{article}
\usepackage{amsmath, amssymb, amsthm}
\usepackage{hyperref}
\usepackage{geometry}
\geometry{margin=1in}

\title{Erd\H{o}s Problem \#724}
\date{Last updated: 2025-08-31}
\author{Source: erdosproblems.com}

\begin{document}
\maketitle

\noindent\textbf{Status:} OPEN \quad \textbf{No prize}

\noindent\textbf{Tags:} combinatorics

\medskip
\noindent\textbf{Problem Statement:}

Let $f(n)$ be the maximum number of mutually orthogonal Latin squares of order $n$. Is it true that\[f(n) \gg n^{1/2}?\]

\bigskip
\noindent\textbf{Remarks:}

Euler conjectured that $f(n)=1$ when $n\equiv 2\pmod{4}$, but this was disproved by Bose, Parker, and Shrikhande \cite{BPS60} who proved $f(n)\geq 2$ for $n\geq 7$.

Chowla, Erd\H{o}s, and Straus \cite{CES60} proved $f(n) \gg n^{1/91}$. Wilson \cite{Wi74} proved $f(n) \gg n^{1/17}$. Beth \cite{Be83c} proved $f(n) \gg n^{1/14.8}$. 

The sequence of $f(n)$ is A001438 in the OEIS.

\bigskip
\noindent\textbf{References:}

[BPS60] Bose, R. C. and Shrikhande, S. S. and Parker, E. T., Further results on the construction of mutually orthogonal
Latin squares and the falsity of Euler's conjecture. Canadian J. Math. (1960), 189-203.
 
 [Be83c] Beth, Thomas, Eine Bemerkung zur Absch\"{a}tzung der Anzahl orthogonaler
lateinischer Quadrate mittels Siebverfahren. Abh. Math. Sem. Univ. Hamburg (1983), 284-288.
 
 [CES60] Chowla, S. and Erd\H{o}s, P. and Straus, E. G., On the maximal number of pairwise orthogonal Latin squares
of a given order. Canadian J. Math. (1960), 204-208.
 
 [Wi74] Wilson, Richard M., Concerning the number of mutually orthogonal Latin squares. Discrete Math. (1974), 181-198.

\bigskip
\noindent\small{Source: \url{https://www.erdosproblems.com/724}}
\end{document}
